\section{A circle family with uniform Fourier bounds}\label{sec:circle}

We first construct representations in which a scalar circle parameter can
later be replaced by a unitary matrix. Two properties are needed: Fourier
coefficients with a summable first moment, uniformly over the group, and
a weak operator limit that retains the parameter.

Let $F$ be the free group on generators $s_j$, $j\geq 1$, with identity
$e$. Put $K=\ell^2(F)$, write $(\delta_x)_{x\in F}$ for its standard
orthonormal basis, and let $p_e$ be the projection onto
$\mathbb C\delta_e$. We use the left regular representation
$l(g)\delta_x=\delta_{gx}$. Throughout the construction,
\[
 t=\frac12,\qquad r=(1-t^2)^{1/2},\qquad
 \mathbb T=\{z\in\mathbb C:|z|=1\}.
\]

\begin{proposition}\label{prop:circle}
There are unitary representations $\pi_z:F\to\B(K)$, indexed by
$z\in\mathbb T$, and operators $A_k(g)\in\B(K)$, indexed by
$k\in\mathbb Z$ and $g\in F$, such that
\begin{align}
 \pi_z(g)&=\sum_{k\in\mathbb Z}z^kA_k(g),
 &\sup_{g\in F}\|A_k(g)\|&\leq c_k,
 &c_k&=\mathbf 1_{\{k=0\}}+\frac{2t^{|k|}}{1-t},
 \label{eq:fourier}\\
 \pi_z(s_j)&\xrightarrow[j\to\infty]{\WOT}tzp_e.
 \label{eq:circle-limit}
\end{align}
The series in \eqref{eq:fourier} converges absolutely in operator norm,
uniformly in $g$ and $z$, and
\begin{equation}\label{eq:moments}
 S_0:=\sum_{k\in\mathbb Z}c_k<\infty,
 \qquad S_1:=\sum_{k\in\mathbb Z}|k|c_k<\infty.
\end{equation}
\end{proposition}

\begin{proof}
The underlying representation is the real-parameter Pytlik--Szwarc
representation \cite[Section~2.3, Equation~(5)]{PS1986}; its cyclic
coefficient is the positive-definite function $g\mapsto t^{|g|}$
of Haagerup \cite[Lemma~1.2]{Haa1979}. We give the tree construction
before introducing the circle parameter.

Use the Cayley tree whose edges join $x$ to $xs_j^{\pm1}$, rooted at
$e$, and write $p(x)$ for the parent of $x\ne e$. Define
\begin{equation}\label{eq:tree-vectors}
 \xi_e=\delta_e,\qquad
 \xi_x=t\xi_{p(x)}+r\delta_x\quad(x\ne e).
\end{equation}
All these vectors have norm one. If $v$ is the last common vertex of the
paths from $e$ to $x$ and $y$, the two tails beyond $v$ are orthogonal
to each other and to $\xi_v$. Thus
\[
 \langle\xi_x,\xi_y\rangle
 =t^{d(x,v)+d(y,v)}=t^{d(x,y)},
\]
where $d$ is the tree distance. Their linear span contains every
$\delta_x$, by \eqref{eq:tree-vectors}. Left multiplication preserves
$d$, so the rule
\[
 \pi(g)\xi_x=\xi_{gx}
\]
extends to a unitary on $K$. These unitaries form a representation.

The difference from the regular representation is confined to a single
finite path. To see this precisely, put $T_g=\pi(g)l(g)^*$.
For $y\ne g$, \eqref{eq:tree-vectors} gives
\begin{equation}\label{eq:root-change}
 T_g\delta_y=
 \frac{\xi_y-t\xi_{g p(g^{-1}y)}}{r},
 \qquad T_g\delta_g=\xi_g.
\end{equation}
The vertex $g p(g^{-1}y)$ is the first step from $y$ toward $g$.
Let $e=h_0,h_1,\ldots,h_d=g$ be the path from $e$ to $g$.
For $y$ outside this path, the first steps toward $e$ and $g$
coincide; hence $T_g\delta_y=\delta_y$. On the path,
\begin{equation}\label{eq:path-action}
 T_g\delta_{h_b}=r\xi_{h_b}-t\delta_{h_{b+1}}\quad(0\leq b<d),
 \qquad T_g\delta_{h_d}=\xi_{h_d}.
\end{equation}
In particular, the path subspace and its orthogonal complement reduce
$T_g$. Expanding the vectors along the path gives
\[
 \xi_{h_b}=t^b\delta_{h_0}
       +r\sum_{a=1}^b t^{b-a}\delta_{h_a}.
\]
By \eqref{eq:path-action}, the subdiagonal entries are $-t$; every
other nonzero entry in row $a$, column $b$ is $t^{|a-b|}$ times
$1$, $r$, or $r^2$. Thus the matrix entries on this subspace satisfy
\begin{equation}\label{eq:path-decay}
 |(T_g)_{a,b}|\leq t^{|a-b|}\qquad(0\leq a,b\leq d).
\end{equation}
This also includes $g=e$, when $T_e=I$.

Let $q:F\to\mathbb Z$ be the homomorphism with $q(s_j)=1$ for every
$j$, and define $D_z\delta_x=z^{-q(x)}\delta_x$. The desired
representation is
\begin{equation}\label{eq:circle-twist}
 \pi_z(g)=z^{q(g)}D_z\pi(g)D_z^*
          =(D_zT_gD_z^*)l(g).
\end{equation}
The first expression proves that it is a unitary representation. The
second follows from $D_zl(g)D_z^*=z^{-q(g)}l(g)$ and supplies its
Fourier bounds.

On the path subspace, the entry in row $a$, column $b$ acquires the
Laurent factor $z^{q(h_b)-q(h_a)}$. Each edge changes $q$ by $1$ or
$-1$, so an entry of exponent $k$ can occur only when $|a-b|\geq|k|$.
For each row or column there are at most two entries at any given
distance $m$. By \eqref{eq:path-decay}, the absolute row and column
sums of the coefficient of $z^k$ are therefore at most
\[
 2\sum_{m\geq|k|}t^m=\frac{2t^{|k|}}{1-t}.
\]
Outside the path the operator is the identity, which contributes only
to the coefficient at $k=0$. The row and column sum bound for operator
norm, followed by multiplication by $l(g)$, proves
\eqref{eq:fourier}. For each fixed $g$ this is a finite Laurent
expansion; the common bounds $c_k$ give the asserted uniform absolute
convergence and \eqref{eq:moments}.

Finally, \eqref{eq:path-action} and \eqref{eq:circle-twist} give
\[
 \pi_z(s_j)\delta_e=tz\delta_e+r\delta_{s_j}.
\]
For any fixed $x\ne e$, all sufficiently large $j$ satisfy
$s_jx\notin\{e,s_j\}$, and then
$\pi_z(s_j)\delta_x=\delta_{s_jx}$. These basis-vector calculations,
together with the uniform norm bound $\|\pi_z(s_j)\|=1$, prove
\eqref{eq:circle-limit}.
\end{proof}
